%%
%% fall-lecture10.tex
%% 
%% Made by Alex Nelson <pqnelson@gmail.com>
%% Login   <alex@lisp>
%% 
%% Started on  2026-10-20T07:34:24-0700
%% Last update 2026-10-20T07:34:24-0700
%% 

\lecture{}

\begin{recall}
We introduced the quantities $\lambda$ and $\mu$ and found
\begin{equation}\label{eq:lec10:stress-tensor}
\sigma_{ij}=-\left[p -\left(\lambda+\frac{2}{3}\mu\right)(\vec{\nabla}\cdot\vec{u})\right]\delta_{i,j}+2\mu\left(S_{ij}-\frac{1}{3}(\vec{\nabla}\cdot\vec{u})\delta_{i,j}\right).
\end{equation}
\end{recall}

\begin{node}
Note that $\lambda$ and $\mu$ are not necessarily constant. In fact,
they are functions of the pressure and temperature
$\lambda=\lambda(p,T)$ and $\mu=\mu(p,T)$.
\end{node}

\begin{node}
We plug this expression for the stress tensor,
Eq~\eqref{eq:lec10:stress-tensor}, into the conservation of linear
momentum, and we find
\begin{subequations}\label{eq:lec10:linear-momentum-conservation}
  \begin{align}
\rho\frac{D u_{i}}{D t} &= -\partial_{i}p + \partial_{j}\tau_{ij}+\gamma_{i}\\
&=-\partial_{i}p+\partial_{j}(\lambda S_{kk})+\partial_{j}(2\mu S_{ij})+\gamma_{i}
  \end{align}
\end{subequations}
where recall our convention has $\tr(S_{ij})=S_{kk}$.
\end{node}

\begin{node}
The conservation of energy $e$ is
\begin{equation}\label{eq:lec10:total-energy-conservation}
\rho\frac{D\varepsilon}{D t}
=-\partial_{i}(pu_{i})+\partial_{i}(\lambda
S_{kk}u_{i})+\partial_{j}(2\mu S_{ij}u_{j})+\gamma_{i}u_{i}
-q_{i}u_{i}+\rho Q.
\end{equation}
\end{node}

\begin{definition}
These two Equations~\eqref{eq:lec10:linear-momentum-conservation} and~\eqref{eq:lec10:linear-momentum-conservation}
plus the conservation of mass equation
are called the \define{Navier--Stokes Equations for Compressible Fluids},
which are 5 equations (1 for mass, 3 for linear momentum, and 1 for energy)
in 6 unknowns ($\rho$, $u_{1}$, $u_{2}$, $u_{3}$, $p$, $e$).
\end{definition}

\begin{definition}
We typically adjoin one extra equation called the \define{Equation of State}
to the Navier--Stokes equations for compressible fluids, which is
usually an equation from thermodynamics.
\end{definition}

\begin{example}
The ideal gas Law is an equation of state
\begin{equation}
p=\rho R T.
\end{equation}
But this  introduces a new unknown, temperature $T$. So we also add to
it the thermodynamic relation
\begin{equation}
e = c_{V}T,
\end{equation}
which gives us two equations of state (technically, I guess).
\end{example}

\begin{example}
We also have the Fourier heat conduction Law
\begin{equation}
q_{i}=-\kappa\partial_{i}T.
\end{equation}
This is a \emph{constitituve law} for heat flux $q_{i}$, since there
are no time derivatives of unknown quantities.
\end{example}

\subsection{Incompressible Navier--Stokes equations}

\begin{remark}
The compressible Navier--Stokes equations are difficult to solve
analytically. Typically, we make the extra assumption that the fluid
is incompressible, which simplifies a lot of difficulties.
\end{remark}

\begin{assume}
The fluid is incompressible: $\rho$ is constant.
\end{assume}

\begin{proposition}
The conservation of mass becomes
\begin{equation}\label{eq:lec10:incompressible:conservation-of-mass}
\partial_{i}u_{i}=0.
\end{equation}
That is to say, the divergence of the velocity vector field vanishes.
\end{proposition}

\begin{node}
The conservation of linear momentum simplifies to
\begin{equation*}\label{eq:lec10:incompressible:conservation-of-linear-momentum}
\rho\frac{D u_{i}}{D t}=-\partial_{i}p+\mu\partial_{j}^{2}u_{i}+\gamma_{i}.
\end{equation*}
\end{node}

\begin{definition}
The conservation of mass
Equation~\eqref{eq:lec10:incompressible:conservation-of-mass} and the
conservation of linear momentum Equations~\eqref{eq:lec10:incompressible:conservation-of-linear-momentum}
are 4 equations in 4 unknowns. These equations are called the
\define{Navier--Stokes Equations for Incompressible Fluids}.
\end{definition}

\begin{node}
The dynamic viscosity $\mu$ is usually a function of pressure and time
$\mu=\mu(p,T)$. However, using the ideal gas Law, for incompressible
flows the dynamic viscosity is a function of temperature alone:
\begin{equation}
\mu=\mu(T).
\end{equation}
So we need to solve
\begin{subequations}
\begin{equation}
e=c_{V}T
\end{equation}
and the conservation of internal energy
\begin{equation}
\rho(\partial_{t}e+\partial_{i}(u_{i}e))=\dots.
\end{equation}
\end{subequations}
\end{node}

\begin{node}
We frequently assume $\mu$ is constant, especially if temperatuer $T$
does not change too much. So frequently the assumption of
incompressibility comes with the assumption of constant viscosity $\mu$.
\end{node}

\begin{note}
The pressure only appears in the conservation of momentum
equations. If we take the divergence of the momentum equations, we
have
\begin{subequations}
\begin{equation}
\partial_{i}[\rho\partial_{t}u_{i}+\rho u_{j}\partial_{j}u_{i}]=-\partial_{i}\partial_{j}p+\partial_{i}[\mu\partial_{j}^{2}u_{i}+\gamma_{i}]
\end{equation}
which can be rewritten as
\begin{equation}
\rho\partial_{t}(\partial_{i}u_{i})+\rho\partial_{i}\partial_{j}(u_{i}u_{j})
=-\partial_{i}\partial_{i}p+\mu\partial_{j}^{2}(\partial_{i}u_{i})+\partial_{i}\gamma_{i}.
\end{equation}
But conservation of mass makes $\partial_{i}u_{i}=0$, simplifying this
to
\begin{equation}
-\rho\partial_{i}\partial_{j}+\partial_{i}\gamma_{i}=\partial_{i}\partial_{i}p=\nabla^{2}p.
\end{equation}
\end{subequations}
This gives us a Poisson equation for pressure. This allows us to
devise an expansion for pressure in terms of velocity, which is very
important in Computational Fluid Dynamics.
\end{note}

\subsection{Special cases of Navier--Stokes equations}

\begin{node}[Incompressible Newtonian fluid with parallel shear flow]
Assume only one component of the fluid velocity is nonzero $u_{1}=u$, $u_{2}=u_{3}=0$.
This simplifies a lot of the math considerably.
\end{node}

\begin{example}[Couette flow]
Recall~(\S\ref{ex:fall:lec09:couette-flow}) the Coeutte flow in
2-dimensions where the top wall had constant velocity $u_{w}$. We saw
how this had velocity $u(y)=u_{w}y/h$. 
\end{example}

\begin{example}
If the wall is ``very far away'' $h\to\infty$, we see that $u(x,y)$ is
a parabola. But conservation of mass gives us
\begin{equation}
\partial_{i}u_{i}=\partial_{x}u(x,y)=0,
\end{equation}
so $u=u(y)$ is a function of $y$ alone, by incompressibility.
\end{example}

\begin{node}
In these situations, the strain rate tensor $S_{ij}$ has only two
nonzero components
\begin{equation}
\S_{ij}=\begin{pmatrix}0 & \frac{1}{2}\partial_{y}u & 0\\
\frac{1}{2}\partial_{y}u & 0 & 0\\
0 & 0 & 0
\end{pmatrix}.
\end{equation}
Similarly, we find
\begin{equation}
\tau_{ij} = \begin{pmatrix}0 & \mu\partial_{y}u & 0\\
\mu\partial_{y}u & 0 & 0\\
0 & 0 & 0
\end{pmatrix}.
\end{equation}
Thus the stress tensor is
\begin{equation}
\sigma_{ij} = \begin{pmatrix}-p & \mu\partial_{y}u & 0\\
\mu\partial_{y}u & -p & 0\\
0 & 0 & -p
\end{pmatrix}.
\end{equation}
\end{node}

\begin{remark}
There are few nice analytic solutions to the Navier--Stokes equations,
and they happen under the parallel shear incompressible flow.
\end{remark}

\subsection{Boundary conditions}

\begin{node}
We have a system of partial differential equations. We need boundary
and initial conditions to solve them.
\end{node}

\begin{definition}
Fluid dynamics uses the term \define{Interface} for boundary between
the fluid and some other medium or material (like a solid wall, or
another fluid when the two fluids are immiscible [like oil and water]).
\end{definition}

\begin{node}[Solid boundary]
We picture the velocity of the fluid matches the velocity of the solid
at the interface. For example, if the wall is stationary, then the
fluid on the surface is stationary. More explicitly: the surface has
outward normal unit vector $\UnitNormalVec{n}$, the surface moves with
velocity $v_{\text{wall}}$, and the velocity vector
field for the fluid $\vec{u}$ has its tangential part equal the
velocity of the wall:
\begin{subequations}
\begin{equation}
  \vec{u}|_{\text{wall}}
  -(\vec{u}|_{\text{wall}}\cdot\UnitNormalVec{n})\UnitNormalVec{n}
  =\vec{v}_{\text{wall}},
\end{equation}
and the fluid's velocity in the ``normal direction of the wall'' vanishes:
\begin{equation}
\vec{u}|_{\text{wall}}\cdot\UnitNormalVec{n}=0.
\end{equation}
\end{subequations}
This is called the \define{No-Slip Boundary Condition}. It was first
observed by Reynolds (but I cannot track the original publication:
people erroneously cite~\cite{reynolds1876force}, but that's clearly
wrong if you read the paper; it is clearly present in a later
publication~\cite[\S20]{reynolds1886theory}, but there may be an
earlier publication).
\end{node}

\begin{node}[Maxwell boundary conditions]
The no-slip boundary condition is admittedly a ``simplifying assumption''.
If we were to ``zoom in'' on the wall, microscopically it would be
``jagged'' and we could try applying geometric optics to determine how
a fluid parcel would scatter off the wall:
\begin{center}
  \includegraphics{img/img.17}
\end{center}
We can model this using a Maxwell-type condition relating the velocity
for the fluid at the wall $\vec{u}|_{\text{wall}}$ to its normal
derivative (in this case, for sake of simplicity, the outward-pointing
normal direction is taken to be the $y$-axis):
\begin{equation}
\vec{u}|_{\text{wall}} =
\frac{2-\beta}{\beta}\lambda\left.\frac{\partial u}{\partial y}\right|_{\text{wall}},
\end{equation}
where $\lambda$ is the mean-free path, $\beta$ is the fraction of
molecules that reflect randomly.

For an idealized smooth wall which
could be engineered by an Ideal Engineer, we would have
$\beta\approx1$. But in this case, we would have
\begin{equation}
\vec{u}|_{\text{wall}} \approx\lambda\left.\frac{\partial u}{\partial y}\right|_{\text{wall}},
\end{equation}
and since $\lambda\ll L$ this reduces to
\begin{equation}
\vec{u}|_{\text{wall}} \approx0.
\end{equation}
Thus we recover the ``no-slip boundary conditions'' as a limiting case
of the Maxwell boundary conditions under a highly idealized setting.
\end{node}

\begin{remark}
Under our assumptions, this Maxwell relation is Reynolds number
independent; if we have water and a hydrophobic surface, then we need
to be more careful.
\end{remark}

\begin{example}
If we have 2 immiscible fluids which do not mix (e.g., oil and water),
then we typically assume:
\begin{enumerate}
\item $\vec{u}_{(1)}=\vec{u}_{(2)}$ at the interface (i.e., in
  tangential directions, and the velocity vectors are the interface
  vanish in the normal directions); \emph{and also}
\item $[\sigma_{ij}n_{j}]_{(1)}=[\sigma_{ij}n_{j}]_{(2)}$ at the
  interface (which intuitively corresponds to ``net force is zero'').
\end{enumerate}
\end{example}

\begin{remark}
In CFD (analytical or theoretical solutions), another type of boundary
conditions is the \define{Inflow/Outflow conditions}. For example, if
we look at a pipe segment, we need to impose some boundary conditions
on the inflow, and we need to impose some boundary conditions on the
outflow. One popular choice is to just impose periodic boundary
conditions, so $\vec{u}|_{\text{inflow}}=\vec{u}|_{\text{outflow}}$.
\end{remark}

\begin{remark}
For more on the boundary conditions,
see Batchelor~\cite[\S1.9]{batchelor1973introduction}.
It is discussed also in White~\cite{white1991viscous} (in \S1-4 and
again briefly in \S2-6), which discusses the Maxwell boundary
conditions among others.

Childress~\cite[\S\S2.6,6.2]{childress2009introduction} offers a
disappointing overview of only the no-slip boundary conditions, but
it's all a Mathematician needs to start solving the equations of motion.
Similarly, Landau and Lifshitz~\cite{landau1987fluid} assert a
boundary condition for ideal fluids in Eq~(2.12) and for viscuous
fluids in Eq~(15.11) are just the no-slip conditions without any
inspection or real
justification. Lamb~\cite[\S79]{lamb1930hydrodynamics} acknowledges
the no-slip boundary condition should be modified when working with
viscous fluids, but it amounts to a stray sentence and little more.

Maxwell proposed his boundary conditions in 1879 according to White~\cite[\S1-4.2]{white1991viscous}.
\end{remark}